\section{Conclusion}

In this paper, we have introduced and analyzed $$-recursive sequence spaces, establishing a novel framework for studying the interplay between computability theory and functional analysis. Our main contribution is the development of a rigorous foundation for examining sequences whose terms are generated by recursive functions, along with their topological and algebraic properties.

The key theoretical result proven in Theorem 1 demonstrates that for any primitive recursive function $f$, the associated sequence space $_f$ is complete under the natural metric $d(x,y) = \sup_{n \in \mathbb{N}} |x_n - y_n|$, providing a crucial structural foundation for future work in this area.

Our approach builds upon recent developments in automated mathematical discovery, particularly the systematic exploration of mathematical structures described in \cite{halachmi2024ramanujan}. However, while previous work has focused primarily on discovering patterns in existing mathematical objects, our framework provides a systematic method for generating entirely new classes of mathematical objects worthy of study.

Several promising directions for future research emerge from this work:
\begin{enumerate}
\item Investigating the relationship between the computational complexity of $f$ and the topological properties of $_f$
\item Extending the theory to partial recursive functions
\item Developing automated methods for discovering and proving properties of $$-recursive sequence spaces
\end{enumerate}

These results lay the groundwork for a broader research program connecting computability theory with classical analysis, potentially leading to new insights in both fields.